\documentclass[10pt]{beamer}

\usetheme[progressbar=frametitle]{metropolis}
\usepackage[italian]{babel}

\usepackage{paralist}

\usepackage{appendixnumberbeamer}

\usepackage{booktabs}
\usepackage[scale=2]{ccicons}

\usepackage{pgfplots}
\usepgfplotslibrary{dateplot}

\usepackage{xspace}
\newcommand{\themename}{\textbf{\textsc{metropolis}}\xspace}

\title{TSP con pick up and delivery}
\subtitle{Progetto del corso di Ricerca Operativa}
\date{13 dicembre 2022}
\author{Michele Vaccari - Matricola 121955}
\institute{Università degli studi di Ferrara\\Corso di laurea magistrale in Ingegneria Informatica e dell'Automazione\\AA 2020-2021}

% logo of my university
\titlegraphic{%
  \begin{picture}(0,0)
    \put(305,0){\makebox(0,0)[rt]{\includegraphics[width=4cm]{../images/logo-unife}}}
  \end{picture}}

\begin{document}

\maketitle

\begin{frame}{Indice}
  \setbeamertemplate{section in toc}[sections numbered]
  \tableofcontents
\end{frame}

\section{Introduzione}

\subsection{Descrizione del problema}
\begin{frame}{\subsecname}

	\metroset{block=fill}
	\begin{block}{Progettino 24 - TSP con pick up and delivery (1 persona)}
		A partire dalla base (nodo $0$ del grafo) un corriere deve soddisfare $n$ richieste di prelievo e consegna di documenti:
		\begin{compactitem}
			\item 	ogni documento è prelevato in un nodo e consegnato in un altro nodo;
			\item ogni nodo è riferito a una singola richiesta, ma nel tragitto tra punto di prelievo e consegna si possono prelevare/consegnare altri documenti.
		\end{compactitem}
		Noto il tempo di percorrenza dei singoli archi, si vuole minimizzare la durata del percorso, con partenza e rientro al deposito.
	\end{block}

\end{frame}

%\begin{frame}[allowframebreaks]{Considerazioni sul problerma}
%	Questo problema di ottimizzazione è NP-hard, poiché coincide con il TSP quando la capacità del veicolo è sufficientemente grande.
%	Anche il problema di verificare se esiste una soluzione ammissibile è un problema fortemente NP-hard, dato che il problema della 3-partizione è un caso particolare (si veda [10] per i dettagli).
%\end{frame}

%\begin{frame}[allowframebreaks]{Interpretazione dei risultati del modello teorico nel problema reale}
%\end{frame}

\begin{frame}[standout]
Grazie per l'attenzione
\end{frame}

\end{document}